\documentclass{article}
\usepackage[UTF8]{ctex}
\usepackage{amsmath,amssymb,amsfonts}
\usepackage{enumitem}
\usepackage{geometry}
\geometry{a4paper, margin=1in}
% === 错题本排版优化（不参与题目标识） ===
\setlength{\parindent}{0pt}
\setlength{\parskip}{0.35em plus 0.10em minus 0.08em}
\setlength{\emergencystretch}{3em}
\setlist{itemsep=0.25em, topsep=0.40em, parsep=0pt}
\allowdisplaybreaks[2]
% === 错题本排版优化结束 ===
\begin{document}
\section*{北京科技大学2025年数学分析真题}
\section*{12．北京科技大学2025年数学分析试卷}

\subsection{题目1}
1．设数列 \(a_{n}=\frac{\sum_{k=0}^{n} \ln C_{n}^{k}}{n^{2}}\) ，证明： \(\lim _{n \rightarrow \infty} a_{n}=\frac{1}{2}\) 。

\subsection{题目2}
2．计算积分 \(\int_{\frac{\pi}{6}}^{\frac{\pi}{3}} \frac{\cos ^{2} x}{x(\pi-2 x)} \mathrm{d} x\) ．

\subsection{题目3}
3．设两个数列 \(\left\{a_{n}\right\}\) 与 \(\left\{b_{n}\right\}\) 的定义如下：

\[
a_{1}=3, b_{1}=2, a_{n+1}=a_{n}+2 b_{n}, b_{n+1}=a_{n}+b_{n}, n=1,2, \cdots .
\]

而数列 \(c_{n}=\frac{a_{n}}{b_{n}}, n=1,2, \cdots\) ，计算极限 \(\lim _{n \rightarrow \infty} c_{n}\) 。

\subsection{题目4}
4．讨论积分 \(\int_{1}^{+\infty} e^{-\alpha x} \frac{\cos x}{x^{p}} \mathrm{~d} x\) 在 \(0 \leq \alpha<+\infty(p>0\) 是常数 \()\) 的一致收敛性．

二．（15 分）若 \(f(x), g(x)\) 在［0，1］上可积，证明： \(\lim _{n \rightarrow \infty} \frac{1}{n} \sum_{i=1}^{n} f\left(\frac{i}{n}\right) g\left(\frac{i-1}{n}\right)=\int_{0}^{1} f(x) g(x) \mathrm{d} x\) ．

三．（15 分）证明不等式：

\begin{itemize}
  \item[（1）] \(\sqrt{t}-\frac{1}{\sqrt{t}}>\ln t(t>1)\) ．
  \item[（2）] \(\sum_{k=1}^{n} \frac{1}{\sqrt{k^{2}+k}}>\ln (1+n)\left(n \in \mathbb{N}^{+}\right)\)．
\end{itemize}

四．（10 分）利用确界原理证明有限覆盖定理．

\begin{itemize}
\end{itemize}

\subsection{题目5}
原题边界“五．”：（15 分）证明：

\begin{itemize}
\item[（1）] 级数 \(\sum_{n=1}^{\infty} \frac{2 x^{2}}{1+n^{3} x^{4}}\) 在 \([0,+\infty)\) 上一致收敛。
  \item[（2）] 级数 \(\sum_{n=1}^{\infty} \frac{n^{2}}{1+n^{4} x^{2}}\) 在 \((0,1)\) 上收敛，但非一致收敛．
\end{itemize}

六．（15 分）若 \(f^{\prime}(x)\) 在 \([a, b]\) 上连续，对任意正整数 \(n\) ，令 \(x_{i}=a+\frac{b-a}{n} i, i=1,2, \cdots, n\) ，记

\[
A_{n}=\frac{b-a}{n} \sum_{i=1}^{n} f\left(x_{i}\right)-\int_{a}^{b} f(x) \mathrm{d} x, n=1,2, \cdots .
\]

证明： \(\lim _{n \rightarrow \infty} n A_{n}=\frac{b-a}{2}[f(b)-f(a)]\) ．

七．（15 分）设函数 \(f(x, y)\) 在单位圆 \(\left\{(x, y) \mid x^{2}+y^{2}=1\right\}\) 上有连续的偏导数，且 \(f(0,1)=f(1,0)\) ，证明：在单位圆上至少有两点满足方程 \(y \frac{\partial}{\partial x} f(x, y)=x \frac{\partial}{\partial y} f(x, y)\) ．

八．（15 分）设 \(\Omega\) 是不含原点的有界闭区域，其体积为 \(V\) ，其边界 \(S\) 为光滑的简单闭曲面， \(\mathbf{n}=\) \((\cos \alpha, \cos \beta, \cos \gamma)\) 为 \(S\) 上点 \(P(x, y, z)\) 处的单位外法向量． \(\mathbf{r}=(x, y, z),(\mathbf{n}, \mathbf{r})\) 为 \(\mathbf{n}\) 与 \(\mathbf{r}\) 的夹角， \(f(t)\) 是 \([0,+\infty)\) 上的连续可微函数，且满足 \(t f^{\prime}(t)+2 f(t)-t=0\) ，计算积分

\[
\iint_{S} f\left(\sqrt{x^{2}+y^{2}+z^{2}}\right) \cos (n, r) \mathrm{d} S
\]

九．（10 分）计算二重积分

\[
\iint_{D}(x \sin \alpha-y \cos \alpha)^{2} \mathrm{~d} x \mathrm{~d} y .
\]

这里 \(D=\left\{(x, y) \left\lvert\, \frac{x^{2}}{a^{2}}+\frac{y^{2}}{b^{2}} \leq 1\right., a>0, b>0\right\}, \alpha\) 为常数．
\end{document}
